\documentclass[11pt,a4paper]{article}
\ifdefined\pdfmapfile
\pdfmapfile{+lm.map}\pdfmapfile{+cm.map}\pdfmapfile{+symbols.map}\pdfmapfile{+euler.map}
\fi
\usepackage[T1]{fontenc}
\usepackage{lmodern,microtype,amsmath,amssymb,amsthm}
\usepackage[left=25mm,right=25mm,top=23mm,bottom=25mm,headheight=14pt]{geometry}
\usepackage{xcolor,fancyhdr}
\definecolor{inkblue}{HTML}{174C60}
\usepackage[unicode,colorlinks=true,linkcolor=inkblue,citecolor=inkblue,urlcolor=inkblue]{hyperref}
\hypersetup{pdftitle={Derivative Newman constants and unconditional results on primes and polynomial zeros}}
\pagestyle{fancy}\fancyhf{}
\fancyhead[L]{\small\itshape Derivative Newman constants and unconditional results}
\fancyhead[R]{\small Research manuscript}\fancyfoot[C]{\thepage}
\renewcommand{\headrulewidth}{0.25pt}
\newtheorem{theorem}{Theorem}[section]
\newtheorem{lemma}[theorem]{Lemma}
\newtheorem{corollary}[theorem]{Corollary}
\newtheorem{proposition}[theorem]{Proposition}
\setcounter{tocdepth}{1}
\setlength{\emergencystretch}{2em}
\raggedbottom
\begin{document}
\thispagestyle{plain}
\begin{center}
{\small\sffamily\color{inkblue} ANALYTIC NUMBER THEORY AND REAL ENTIRE FUNCTIONS}\\[10pt]
{\LARGE\bfseries Derivative Newman constants and unconditional}\\[4pt]
{\LARGE\bfseries results on primes and polynomial zeros}\\[9pt]
{\small Research manuscript \qquad 21 November 2026}
\end{center}
\begin{abstract}
The derivative de Bruijn--Newman constants of Ki, Kim and Lee are
nonnegative in every order and decrease to zero. Under the ordinary
Riemann hypothesis they are all zero. The negative-time proof transfers
Dobner's approximation to fixed linear differential operators and yields
a positive lower bound for the number of distinct zeros per unit height
in a logarithmically displaced strip. Littlewood's oscillation theorem
refutes the 2010 monotonicity conjecture of Choie, Planat and Sol\'e
without a prime-gap hypothesis. Both associated primorial ratios have
arbitrarily long common runs in each direction. We also obtain
unbounded signed oscillations on the primes in the Zhang race without
assuming RH. For factorially weighted polynomial multipliers we
characterize exactly which polynomials with nonnegative coefficients
give complex zero decreasing sequences, settling a question of
Csordas and Forg\'acs. We give
an exact degree-72 counterexample for their specified quadratic.
For permuted gaps we prove Tom\'as Garc\'ia's proposed sharp upper
constant and quantify, without RH, how rarely a permutation of Farey
gaps has discrepancy as small as the native order.
\end{abstract}
\tableofcontents
\medskip
\noindent Throughout, RH concerns the ordinary Riemann zeta function.
Theorems in this manuscript are proposed with proofs for independent
verification. Bibliographic comparisons identify earlier results and
their hypotheses; they do not certify historical priority.
\par\smallskip
\noindent The motivating corpus includes a treatment of zero collisions
under heat flow (item 126) and an exact sign certificate for a multiplier
sequence (item 187). The arguments below use the published analytic inputs
listed in the references and an independent exact certificate for the
factorially weighted sequence. \href{https://github.com/QEEDDD/Verifiable-Mathematical-Reasoning-Corpus/tree/c2ac34ea4622778187173f41bb1996deee8081b7}{Corpus snapshot}.
\newpage

\section{The derivative de Bruijn--Newman constants}\label{sec:derivativeflow}

Ki, Kim and Lee defined the derivative de Bruijn--Newman constants in
2009 \cite{KKL,LM}. We use classical time $t$ with
$H_t(z)=\int_0^\infty e^{tu^2}\Phi(u)\cos(zu)\,du$ and denote
by $\Lambda^{(m)}$ the infimum of the times for which every zero of
$H_t^{(m)}$ is real. Here
\[
 \Phi(u)=\sum_{n\ge1}(2\pi^2 n^4e^{9u}-3\pi n^2e^{5u})
                          \exp(-\pi n^2e^{4u}).
\]
The convention is $\Lambda^{(m)}=4\lambda_m$
of Lagarias--Montague, equations (1.7)--(1.8), because their integral
uses the scale $u\mapsto 2u$. In particular $\Lambda^{(0)}=\Lambda$.

\medskip\noindent
\textbf{Theorem.} \emph{Fix an integer $m\ge0$ and a real $t<0$.
The function $H_t^{(m)}$ has infinitely many nonreal zeros.
In particular $\Lambda^{(m)}\ge0$ for every $m$, and unconditionally}
\[
 \lim_{m\to\infty}\Lambda^{(m)}=0.
\]
\emph{If the ordinary RH holds, then $\Lambda^{(m)}=0$ simultaneously
for every $m\ge0$.}

\medskip\noindent
\emph{Proof.} Use Dobner's normalization
$\xi_t((1+iz)/2)=8H_t(z)$. This identity follows from his
equations (1)--(2) and evenness of the kernel. Thus the derivative
in the complex variable of $\xi_t$ corresponds to a nonzero
constant multiple of a spatial derivative of $H_t$.
For fixed $t<0$ put
\[
 \begin{split}
 Z_t(s)&=\sum_{n\ge1}e^{t(\log n)^2/4}n^{-s},\qquad
 J(s)=s+\frac{|t|}{4}\operatorname{Log}\frac{s}{2\pi},\\
 \gamma(s)&=\tfrac12s(s-1)\pi^{-s/2}\Gamma(s/2),\qquad
 B(s)=\gamma(s)\exp\left(\frac{(s-J(s))^2}{|t|}\right).
 \end{split}
\]
Use the principal logarithm high in the upper half-plane. Dobner's Theorem 4 \cite{Dobner}, specialized to $\zeta$, gives
\[
 \frac{\xi_t(J(s))}{B(s)}=Z_t(s)+o(1)
 \qquad (\Im s\to+\infty)
\]
uniformly on each fixed strip $a\le\Re s\le b$. Here $B$ is holomorphic
and nonzero at sufficiently large heights. Dobner's Lemma 3 shows
that $Z_t$ has a zero. We need an approximation to the spatial
derivatives with the same leading Dirichlet series.

Set $s=w+iY$, with $w$ ranging over a fixed compact subset of
$\mathbb C$ and $Y\to\infty$. For $Y>2\pi$ define the real constant
$c_Y=\tfrac12\log(Y/(2\pi))$, and put
\[
 A_{m,Y}(w)=\frac{\xi_t^{(m)}(J(s))}{B(s)c_Y^m}.
\]
We claim, for every fixed $m$, that
\[
 A_{m,Y}(w)=Z_t(w+iY)+o(1)
\]
locally uniformly in $w$. Dobner gives $m=0$. Stirling's formula
for the digamma function in sectors separated from the negative
real axis gives, uniformly for bounded $w$,
\[
 J'(s)=1+O(Y^{-1}),\qquad
 \frac{B'(s)}{B(s)}
  =\tfrac12\operatorname{Log}\frac{s}{2\pi}
       +O_t\!\left(\frac{\log Y}{Y}\right).
\]
Indeed the logarithmic derivative of $s(s-1)$ is $O(Y^{-1})$,
and differentiating the squared logarithm in the exponential of $B$
contributes $O_t((\log Y)/Y)$. Hence $(B'/B)/c_Y\to1$.
Suppose the claim holds at $m$. The series for $Z_t$ is absolutely
bounded on each fixed vertical strip, so $A_{m,Y}$ is locally uniformly
bounded even on a fixed enlargement of the compact set.
Cauchy's estimate gives $\partial_w A_{m,Y}=O(1)$ on the original
compact set. The exact recursion is
\[
 A_{m+1,Y}(w)=\frac{1}{J'(s)}
 \left(\frac{\partial_w A_{m,Y}(w)}{c_Y}+
       \frac{B'(s)}{c_Y B(s)}A_{m,Y}(w)\right).
\]
This proves the claim by induction without differentiating a
pointwise error estimate.

Choose a closed disk $D$ about a zero of $Z_t$ with no zero on
$\partial D$ and set $\delta=\min_{\partial D}|Z_t|>0$.
Vertical almost periodicity (Dobner's Theorem 5) supplies arbitrarily
large $Y$ with
$\sup_{w\in\partial D}|Z_t(w+iY)-Z_t(w)|<\delta/3$.
One may also obtain these shifts directly by truncating the Dirichlet
series uniformly on the strip of $D$ and simultaneously returning
the finitely many phases $e^{-iY\log n}$ close to one. On these
translates the proved approximation gives
$|A_{m,Y}-Z_t|<2\delta/3$ on $\partial D$ for sufficiently large $Y$.
Rouch\'e's theorem yields $w_Y\in D$ with
$\xi_t^{(m)}(J(w_Y+iY))=0$. For such $w_Y$,
\[
 \Re J(w_Y+iY)=\tfrac{|t|}{4}\log Y+O_t(1)\to+\infty,
 \qquad \Im J(w_Y+iY)=Y+O_t(1)\to+\infty .
\]
The zeros have unbounded heights and eventually lie to the right of
the critical line. They yield the asserted nonreal zeros after the
linear change of variable.

It follows that every $\Lambda^{(m)}$ is nonnegative.
Ki--Kim--Lee proved
$\Lambda^{(m+1)}\le\Lambda^{(m)}$ and
$\lim_{m\to\infty}\Lambda^{(m)}\le0$. Combining these results gives
the displayed unconditional limit. Under RH the order-zero function
$H_0$ belongs to the Laguerre--P\'olya class, so every $H_0^{(m)}$
has only real zeros. Time zero therefore realizes every infimum.\hfill$\square$

\medskip\noindent
This uses the constants originally defined by Ki--Kim--Lee. Their
published monotonicity and upper bound on the limit are inputs.
Dobner states the approximation and the zeros conclusion for the
flow before spatial differentiation. The differentiation argument gives the same negative-time
obstruction at each fixed derivative order. Ki's 2022 report records the upper bound for this limit
\cite{Ki22}. We have not located a separately named conjecture stating
the equality for derivatives. The order-zero lower bound is the
Rodgers--Tao theorem \cite{RT}.

\subsection{Fixed differential operators and the number of zeros}
Let $P(X)=a_dX^d+\cdots+a_0$ be any nonzero complex polynomial and
let $D_s$ be differentiation in the xi variable.  For each fixed $t<0$
there are $L_t,C_t>0$ such that every sufficiently high interval of
heights of length $L_t$ contains a zero of $P(D_s)\xi_t$ in
\[
 \left|\Re s-\tfrac{|t|}{4}\log(\Im s)\right|\le C_t .
\]
The initial height can depend on $P$. In particular the number of
distinct such zeros at heights at most $T$ is $\gg_t T$ for large $T$.

To prove this, divide $P(D_s)\xi_t(J(w+iY))$ by $a_d B(w+iY)c_Y^d$.
The preceding induction gives $Z_t(w+iY)+o(1)$ locally uniformly, since
each derivative of smaller order carries a factor $c_Y^{j-d}$.
Fix the disk and its boundary tolerance as above, before choosing $P$.
The permitted return heights are relatively dense.
Indeed, truncate the absolutely convergent series uniformly on the disk.
The closure of the finite phase orbit
$(e^{-iY\log k})_{k\le K}$ is a compact group.
Finitely many translates of any sufficiently small identity
neighborhood by orbit elements cover that group. If $a_1,\ldots,a_b$
are their translation times, for every $Y$ at least one
$Y-a_i$ returns to the neighborhood. The truncation tail then gives
the required tolerance for $Z_t$. The imaginary coordinate of the resulting
zero differs from its return height by a bounded amount, and the formula
for $\Re J$ gives the indicated strip. Enlarging the interval length
absorbs both bounded shifts. Selecting alternate intervals gives
distinct zeros.

When $P$ has only real roots, one can also apply $P(D_z)$ to the
classical $H_t$. Define $\Lambda_P$ as the infimum of the times at which
all the zeros of $P(D_z)H_t$ are real. The affine change of variables
in the theorem and the preceding extension exclude every $t<0$.
At $t\ge\Lambda$, $H_t$ is a real entire function of order at most one
with only real zeros, hence is in the Laguerre--P\'olya class.
Each factor $D_z-a$ with $a\in\mathbb R$ preserves this class:
on real-rooted polynomials the logarithmic derivative is a sum of
simple real fractions with positive residues, so it takes a real value
$a$ only on the real axis away from the zeros; multiple zeros are
handled by factoring. Locally uniform polynomial approximation and
Hurwitz's theorem extend this to the class. The result of the operator is nonzero by uniqueness of the Fourier
transform, since the integrable Fourier kernel is nonzero on an interval
and is multiplied by a nonzero polynomial.
Consequently $0\le\Lambda_P\le\Lambda$ and under RH $\Lambda_P=0$.

\section{An unconditional disproof of Choie--Planat--Sol\'e Conjecture 1}
Let $p_n$ denote the $n$th prime, $N_n=\prod_{j\le n}p_j$,
$\vartheta(x)=\sum_{p\le x}\log p$,
and $\Psi(m)=m\prod_{p\mid m}(1+1/p)$. Consider the original ratios
\[
 u_n=\frac{N_n}{\varphi(N_n)\log\log N_n},\qquad
 v_n=\frac{\Psi(N_n)}{N_n\log\log N_n}.
\]
In \emph{On Nicolas criterion for the Riemann Hypothesis}
(2010, arXiv:1012.3613v2), Conjecture~1 of Choie, Planat and Sol\'e
\cite{cps} asserts the adjacent inequality (1) of their Proposition~3 for every
$n\ge1$. For $n\ge2$ this is equivalent to $u_n>u_{n+1}$.
At $n=1$ the denominator $\log\log 2$ is negative, so one cannot
cross-multiply there as in their Proposition~3. All violations below
occur at arbitrarily large indices of the printed inequality itself.
Their Propositions~4 and~7 give infinitely many violations for the two
ratios under Cram\'er's conjecture on prime gaps.
Their Propositions~5 and~8 treat the opposite direction unconditionally.

\medskip\noindent\textbf{Theorem.}
Both sequences have arbitrarily long simultaneous runs of strict increase
and arbitrarily long simultaneous runs of strict decrease, unconditionally.
Each direction has runs of length at least
$c_1\sqrt{q}\log\log\log q/\log q$ with largest entering prime
$q\to\infty$, for an absolute $c_1>0$. In particular, Conjecture~1
is false at infinitely many indices.

\medskip\noindent\emph{Proof.}
At a step entering the prime $q=p_j$, write $A=\vartheta(q)-\log q>1$
and $d=\log(1+\log q/A)$. Exact division of consecutive terms gives
\[
 \frac{v_j}{v_{j-1}}=
 \frac{1+1/q}{1+d/\log A},\qquad
 \frac{u_j}{u_{j-1}}=
 \frac{1+1/(q-1)}{1+d/\log A}.
\]
When $A\ge q$, we have
$d<\log q/A\le \log A/q$, so both ratios rise.
When $A+\log q\le q-1$, the inequality $\log(1+y)>y/(1+y)$ gives
$d>\log q/(A+\log q)>\log A/(q-1)$, so both fall.

By Littlewood's oscillation theorem in its Chebyshev form
(Ingham, \emph{The Distribution of Prime Numbers}, Thm.~6.3; see also \cite[(2)--(3)]{PT}), there
is a constant $c>0$ for which, on separate sequences of $x\to\infty$,
\[
 \vartheta(x)-x>E(x)
 \quad\hbox{and}\quad
 \vartheta(x)-x<-E(x),\qquad
 E(x)=c\sqrt{x}\log\log\log x.
\]
Take $L=\lfloor E(x)/(3\log(2x))\rfloor$.
In the positive case take the last $L$ primes at or below $x$. There
are that many by the prime number theorem. At each of these primes
$q$, subtracting at most $L$ log-prime increments gives
$A\ge\vartheta(x)-L\log x>x\ge q$. All $L$ steps rise.
In the negative case take the first $L$ primes above $x$. The prime
number theorem puts them all in $(x,2x]$ for large $x$.
At every such prime,
$A+\log q=\vartheta(q)\le\vartheta(x)+L\log(2x)
<x-2E(x)/3<q-1$. All $L$ steps fall. Since $L\to\infty$,
this proves the existence of arbitrarily long runs. The largest
entering prime is in $[x/2,2x]$ for large $x$, by the prime number
theorem. Replacing $x$ by this prime in the order of magnitude for $L$
gives the stated quantitative bound.\hfill$\square$

\medskip\noindent\textbf{An explicit bound.}
Platt and Trudgian, \emph{On the first sign change of
$\vartheta(x)-x$}, \emph{Mathematics of Computation} \textbf{85}
(2016), 1539--1547, p.~1545 \cite{PT}, prove the existence of an
$x<\exp(727.951332668)$ with $\vartheta(x)-x>10^{152}$.
There are more than $10^{149}$ primes up to that $x$, since
$\pi(x)\log x\ge\vartheta(x)>10^{152}$ and $\log x<728$.
For each of the last $10^{149}$ primes the proof above gives
$A\ge\vartheta(x)-728\cdot10^{149}>x$.
Thus at least $10^{149}$ consecutive steps of both original ratios
strictly increase before an entering prime reaches
$\exp(727.951332668)$.

The positive case with $L=1$ explains how the original gap hypothesis
disappears. At a large positive excursion, use the prime step
\emph{ending} at the witnessing prime. Removing that prime from
$\vartheta$ costs its logarithm, already smaller than the certified
excursion. The conclusion concerns adjacent primorial ratios.
The separate inequality $u_n>e^\gamma$ at every sufficiently large
index, central to Nicolas's RH criterion, is not decided by these runs.

Karttunen also recorded the inequality
$\vartheta(q)^q> (\vartheta(q)-\log q)^{q+1}$ at every entering prime as
a conjectural observation in OEIS A342000 \cite{oeis1}. For $A>1$
this says $q d>\log A$, precisely the strict decrease condition for
$v_j$. The increasing runs refute that observation as well.

\subsection{A question about fifth powers of primorials}
Karttunen's OEIS entry A342455 asks whether the increase observed after the fifth
term continues for every subsequent term \cite{oeis5}.  For each integer $k\ge1$,
write
\[
 W_k(n)=\frac{\sigma(N_n^k)}
 {e^\gamma N_n^k\log\log(N_n^k)} .
\]

\begin{theorem}\label{thm:fifth}
For every $k\ge5$ and $n\ge5$ one has $W_k(n+1)>W_k(n)$.
In particular the answer to the question in A342455 is yes.
\end{theorem}
\begin{proof}
Put $q=p_{n+1}$, $s=\log q$, $A=\vartheta(q)-s=\log N_n$, and
$g_k(q)=\sum_{j=1}^k q^{-j}$.
The new Euler factor in $\sigma(N_n^k)/N_n^k$ is $1+g_k(q)$,
while the logarithmic denominator changes from $\log(kA)$ to
$\log(k(A+s))$.  Both are positive for $n\ge5$.  Consequently a strict rise
is equivalent to
\[
 D_k(q):=g_k(q)\log(kA)-\log(1+s/A)>0.
\]
In this range $g_k(q)$ and $\log(kA)$ are positive and increasing in $k$,
so it suffices to check $k=5$.
Since $g_5(q)>1/q$ and $\log(1+s/A)<s/A$,
the inequality $A\log(5A)>qs$ suffices.
For $q\ge149$ we have $s>5$.  Letting $x$ increase to $q$ in
the estimate $\vartheta(x)>x(1-1/\log x)$ for $x\ge41$
\cite[(3.16)]{rs62} gives $A\ge q(1-1/s)\ge4q/5$.  Hence
\[
 \frac{A\log(5A)}q\ge(1-1/s)(s+\log4)
 \ge s-1+\tfrac45\log4>s.
\]
For the 29 primes $13\le q\le139$ the following outward-rounded
interval certificate gives $D_5(q)>1/625$ at every step.
The loop updates $A$ only after testing the step, so the prime $q$ is
excluded from $A$.  The code uses SageMath syntax.
\begin{verbatim}
R = RealIntervalField(128)
A = R(0)
for q in prime_range(2, 149):
    s = R(q).log()
    if q >= 13:
        g = sum(R(q)^(-j) for j in (1,2,3,4,5))
        d = g*(5*A).log() - (1+s/A).log()
        assert d.lower() > QQ(1)/625
    A += s
\end{verbatim}
This covers every right endpoint $q=p_{n+1}$ with $n\ge5$.
\end{proof}

\section{Unconditional oscillation of the Zhang prime race}
Following Lichtman and Pomerance \cite[Sec.~5]{lpzh}, put
\[
 F(x)=\frac1{\log x}-\sum_{p\ge x}\frac1{p\log p},
 \qquad E_Z(x)=\sqrt{x}\log^2x\,F(x)\quad(x>1).
\]
A prime $p$ is Zhang if $F(p)\ge0$.  The relative logarithmic density
predicted for these primes was obtained under RH and linear independence
by Lichtman, Martin and Pomerance \cite[Thm.~5.6]{lmp}.
The following unconditional assertion concerns the two signs.
\begin{theorem}\label{thm:zhang}
Along the primes, $\limsup E_Z(p)=+\infty$ and
$\liminf E_Z(p)=-\infty$.  In particular there are infinitely many
strict Zhang primes and infinitely many non-Zhang primes.
\end{theorem}
\begin{proof}
First suppose RH. For the Mertens product and its error use
\[
P(x)=\prod_{p\le x}(1-1/p)^{-1},\qquad
E_M(x)=\sqrt{x}\log x\,\log\bigl(P(x)/(e^\gamma\log x)\bigr).
\]
The explicit-formula comparison \cite[(5.12)]{lmp} gives
$E_Z(x)=E_M(x)+O(1/\log x)$ under RH.  Although it appears within
a proof of a density result stated under linear independence as well,
(5.12) itself follows by comparing the RH explicit formulas cited there.
Mertens' product theorem and $\log(1+y)/y\to1$ give
\[
 E_M(x)= (e^{-\gamma}+o(1))\sqrt{x}\{P(x)-e^\gamma\log x\}.
\]
Diamond and Pintz \cite[Thm.~1.1]{dp} prove that the expression on
the right without its positive factor is unbounded in both directions.

If RH is false, choose a zero $\rho$ with $\beta=\Re\rho>1/2$, and
fix $1-\beta<b<1/2$.  We show that for arbitrarily large $x$ each of
$F(x)>x^{-b}$ and $F(x)<-x^{-b}$ holds. Take $a=3/2$ and
$H(s)=\int_a^\infty F(x)x^{-s-1}\,dx$ for $\Re s>0$.
The distributional derivative of $F$ consists of the density
$-1/(x\log^2x)$ and the point masses $1/(p\log p)$ at primes.
Integration by parts followed by two derivatives in $s$ therefore gives
\[
 \big(sH(s)-a^{-s}F(a)\big)''
 =-\frac{\zeta'}{\zeta}(1+s)
 -\sum_{p,\,j\ge2}\frac{\log p}{p^{j(1+s)}}-\frac{a^{-s}}s.
\]
The prime-power series converges locally absolutely for $\Re s>-1/2$.
The two simple poles at zero cancel, and $\zeta$ has no real zeros in
$(0,1)$. Integrating twice along a neighborhood of the real axis thus
continues $sH(s)$ analytically past every real $s>-1/2$.
Since $F(x)\to0$, $sH(s)\to0$ as $s\downarrow0$, so $H$ is also
regular at zero. At $s=\rho-1$ the displayed second derivative has
a genuine pole, precluding holomorphy of $H$ on any half-plane
containing that point.

For either sign, the transform of $F(x)\pm x^{-b}$ is
$H(s)\pm a^{-s-b}/(s+b)$. It is regular along the real axis to the
right of $-b$ and cannot be holomorphic on any half-plane containing
$\rho-1$. If $F(x)\pm x^{-b}$ were eventually of one sign,
Landau's oscillation theorem for Mellin transforms
\cite[Sec.~3]{dp} would give a real singularity at its abscissa of
convergence, which is at least $\beta-1>-b$.  This contradiction
proves the claim. Since $b<1/2$, the claim implies both asserted
divergences of $E_Z(x)$ along real arguments.

Finally, for successive primes $p<x\le q$, the prime tail in $F$
is constant on $(p,q]$, so $F(q)\le F(x)$ and
$F(p)\ge F(x)-1/(p\log p)$. The prime number theorem gives
$p\sim x\sim q$, and $\sqrt{p}\log^2p/(p\log p)\to0$.
Thus both divergences pass to prime arguments. Points $x$ that are
primes already need no transfer.
\end{proof}

\section{Polynomial weights on factorial multiplier sequences}
\label{sec:czds}

For a real sequence $\gamma$ write $T_\gamma(\sum c_n x^n)=\sum
\gamma_n c_n x^n$. Let $Z_c$ and $Z_{\mathbb R}$ count nonreal and
real roots of a nonzero polynomial, with multiplicity. A sequence is
\emph{complex zero decreasing} (CZDS) if
$Z_c(T_\gamma p)\le Z_c(p)$ for all real polynomials $p$, using the
usual convention $Z_c(0)=0$.

Csordas and Forg\'acs ask for conditions on a polynomial $p$ with
nonnegative coefficients for $(p(n)/n!)$ to be a CZDS. They single out
$p(x)=x^2+x+1$ as an open question and prove that the corresponding
sequence preserves real-rootedness \cite[Example 34 and Problems 36(b),
37 in arXiv v2]{CFbessel}. We obtain an if-and-only-if condition.

\begin{theorem}\label{thm:czdsfamily}
Let $d>0$. For each nonzero polynomial $p\in\mathbb R[x]$ with
nonnegative coefficients, the sequence $(p(n)/\Gamma(n+d))_{n\ge0}$
is complex zero decreasing if and only if every zero of $p$ is real.
In particular, $( (n^2+n+1)/n!)_{n\ge0}$ is not complex zero decreasing.
\end{theorem}

Bakan and Ruscheweyh proved the converse to Laguerre's interpolation
theorem under the Stieltjes Carleman condition
$\sum_{n\ge1}m_n^{-1/(2n)}=\infty$
\cite[Theorem 1, Corollary 1(B), and Sections 2--4]{BRczds}.
Their proof passes from weighted polynomial approximation for dilates
of a positive representing measure to a variation-diminishing
convolution. We give the signed-measure version needed to identify the
interpolant for each sequence considered here. For a real function $f$ on an ordered set, let
$V(f)$ be the supremum, over all finite increasing sequences of
arguments, of the number of sign changes after zero values are removed.

\begin{lemma}\label{lem:czdsmellin}
Let $\nu$ be a finite real signed measure on $(0,\infty)$ such that
$\int e^{c\sqrt{u}}\,d|\nu|(u)<\infty$ for some $c>0$. Suppose
$m_n=\int u^n\,d\nu(u)>0$ for $n\ge0$. Suppose also that
$M(z)=\int u^z\,d\nu(u)$ converges absolutely on an open vertical
strip containing zero. If $(1/m_n)$ is a CZDS, the reciprocal of
$M$ near zero extends to an entire function with only real zeros.
\end{lemma}

\begin{proof}
The positive entries $m_n$ ensure preservation of degree. Applying the
CZDS inequality to $T_m h$ and subtracting from the degree gives
\[
   Z_{\mathbb R}(T_m h)\le Z_{\mathbb R}(h)
         \qquad(0\ne h\in\mathbb R[x]).
\]
The operator $T_m$ on polynomials is integration against dilates:
$(T_m h)(x)=\int h(xu)\,d\nu(u)$. We prove the variation bound
\begin{equation}\label{eq:czdsvar}
    V\!\left(x\mapsto\int g(xu)\,d\nu(u)\right)\le V(g)
\end{equation}
on $(0,\infty)$ for every bounded continuous real $g$.

It suffices to consider nonzero $g$ with $V(g)=k<\infty$.
The nonzero values of $g$ form $k+1$ successive sign groups.
Choose a point where $g=0$ between each two consecutive groups.
A polynomial $Q$ with these $k$ points as simple roots can be given a
leading sign such that $g/Q\ge0$ wherever $Q\ne0$. For $k=0$, use
$Q=1$ or $Q=-1$. Fix positive evaluation points $x_1<\cdots<x_l$,
write $D_x(u)=xu$, and introduce the finite positive measure
\[
     \tau(dy)=|Q(y)|\sum_{i=1}^l (D_{x_i})_*|\nu|(dy).
\]
Polynomial growth and the assumption on $|\nu|$ give
$\int e^{c'\sqrt y}\,d\tau(y)<\infty$ for a sufficiently small
$c'>0$.

Real polynomials are dense in $L^2(\tau)$. To see this, lift $\tau$
symmetrically by mapping $y$ to $\sqrt y$ and $-\sqrt y$ with equal
weights. The resulting measure $\sigma$ on $\mathbb R$ has
an exponential moment. If $H\in L^2(\sigma)$ is orthogonal to every
polynomial, $\int e^{zt}H(t)\,d\sigma(t)$ is holomorphic for
$|\operatorname{Re}z|<c'/2$ by Cauchy--Schwarz. All its derivatives
at zero vanish. Fourier uniqueness gives $H=0$ almost everywhere.
Even functions of $t$ can thus be approximated by even polynomials,
by symmetrizing the approximants. Pulling these back by $y=t^2$ proves
the asserted density.

Put $h(y)=\sqrt{g(y)/Q(y)}$ when $Q(y)\ne0$ and set it equal to zero
at the roots. Then $g=Qh^2$ everywhere and $h\in L^2(\tau)$
because $g$ is bounded. Take real polynomials $P_j\to h$ in
$L^2(\tau)$. Cauchy--Schwarz gives convergence of $P_j^2$ to $h^2$
in $L^1(\tau)$. Consequently for each $i$,
\[
 T_m\!\left(Q(P_j^2+1/j)\right)(x_i)
        \longrightarrow\int g(x_i u)\,d\nu(u).
\]
Every polynomial on the left has at most $k$ real zeros, since the
input has precisely the simple real roots of $Q$. If the limiting
values at the selected points changed sign more than $k$ times, a
subsequence of $k+2$ nonzero alternating values would retain its signs
for sufficiently large $j$, contradicting the intermediate value
theorem. This proves \eqref{eq:czdsvar}.

Push $\nu$ forward under $u\mapsto-\log u$ and denote the resulting
measure by $\lambda$. Substituting $g(x)=F(\log x)$ in
\eqref{eq:czdsvar} gives $V(F*\lambda)\le V(F)$ for every bounded
continuous $F$ on $\mathbb R$. Schoenberg's characterization of
variation-diminishing convolutions \cite[Theorem 2, p.~99]{ISczds}
states that the bilateral Laplace transform of such a nonzero measure
is the reciprocal of a Laguerre--P\'olya entire function in a strip
about the imaginary axis, up to a nonzero real constant. The transform
here is $\int e^{-zt}\,d\lambda(t)=M(z)$, and $M(0)=m_0>0$.
\end{proof}

\begin{proof}[Proof of Theorem~\ref{thm:czdsfamily}]
If all zeros of $p$ are real, nonnegative coefficients imply they are
nonpositive. Then $p(z)/\Gamma(z+d)$ is a real Laguerre--P\'olya
function with only nonpositive zeros, by the Weierstrass product for
$1/\Gamma$. Laguerre's complex zero decreasing theorem gives sufficiency
\cite[Theorem 4.1(3)]{CCsurveyZero}.

Conversely assume the zero inequality holds. For a constant polynomial
the necessity is vacuous. Otherwise write
$p(z)=A\prod_{j=1}^r(z-\alpha_j)$ with $A>0$ and choose an integer
$k>\max(0,\max_j\operatorname{Re}\alpha_j)$. If $\gamma$ is a CZDS
then so is $(\gamma_{n+k})_{n\ge0}$, since
$T_\gamma(x^kh)=x^kT_{(\gamma_{n+k})}h$ and the extra roots are real.

Take the product of the positive measure
$s^{k+d-1}e^{-s}\,ds$ on $(0,\infty)$ and the complex measures
$u_j^{k-\alpha_j-1}\,du_j$ on $(0,1)$ for $1\le j\le r$.
Divide by $A$ and push the product forward under
$(s,u_1,\ldots,u_r)\mapsto s\prod_j u_j$. This defines a finite
real signed measure $\nu$: conjugating permutes the factors indexed
by conjugate roots. Its total variation has an exponential moment
in $u$ itself, since $s\prod_j u_j\le s$. Integration of the factors
for $\operatorname{Re}z$ close to zero gives
\[
    M(z)=\frac{\Gamma(z+k+d)}{p(z+k)}.
\]
Its moments at nonnegative integers are reciprocals of the shifted CZDS.
Lemma~\ref{lem:czdsmellin} shows that the entire function
$p(z+k)/\Gamma(z+k+d)$ has only real zeros. A nonreal zero of $p$
would persist in this product. This proves necessity. For $d=1$ and
$p(x)=x^2+x+1$ the roots are $(-1\pm i\sqrt3)/2$.
\end{proof}

\begin{theorem}\label{thm:czdssqrt}
The sequence $(\sqrt{n}/n!)_{n\ge0}$ fails to be complex zero
decreasing, answering Csordas--Forg\'acs Problem 39(b) in arXiv v2.
\end{theorem}

\begin{proof}
If it were a CZDS, shifting once would give a positive CZDS with
reciprocal terms $(n+1)!/\sqrt{n+1}$. A positive measure with these
moments is the law, up to normalization, obtained by multiplying
independent variables with densities $s e^{-s}$ on $(0,\infty)$ and
$1/\sqrt{\pi\log(1/u)}$ on $(0,1)$. Substitution $v=\log(1/u)$
shows its Mellin transform is $M(z)=\Gamma(z+2)/\sqrt{z+1}$
for $\operatorname{Re}z>-1$, taking the branch positive for real $z>-1$.
The product is bounded by $s$, so the integrability hypothesis of
Lemma~\ref{lem:czdsmellin} holds. The resulting entire extension $E$ of
$1/M$ would satisfy $E(z)^2=(z+1)/\Gamma(z+2)^2$ everywhere, by
analytic continuation. The right side has a simple zero at $-1$,
whereas every zero of a square has even multiplicity.
\end{proof}

Craven and Csordas had previously ruled out the CZDS property for
$(1+n+n^2)$ via a reciprocal Hankel matrix
\cite[Example 4.5]{CCsurveyZero}, and for $((n-1/2)^2/n!)$ by a finite
determinant \cite[Theorem 3.3]{CCserdica1996}. The Carleman condition of Bakan and Ruscheweyh holds in each application
above by Stirling's formula, and their general criterion supplies the
interpolation principle. The explicit multiplicative measures determine
the interpolating function and the zeros or branch point obstructing it.
Theorem~\ref{thm:czdsfamily} characterizes the factorial family across
all polynomial degrees.

\begin{theorem}\label{thm:czdstree}
The sequence $\bigl(1/[n!(n+1)^{n+1}]\bigr)_{n\ge0}$ fails to be
complex zero decreasing, answering Csordas--Forg\'acs Problem 43(b)
in arXiv v2.
\end{theorem}

\begin{proof}
We give a construction of a positive random variable $S$ with
$\mathbb E S^z=(1+z)^{1+z}$ for $\operatorname{Re}z>-1$,
using the logarithm real on $(0,\infty)$. Let $N$ be a Poisson random
measure on $(0,\infty)$ with intensity
$\eta(dt)=t^{-2}e^{-t}\,dt$. With $\epsilon>0$, define
\[
       Y_\epsilon=1-\int_\epsilon^\infty t\,(N-\eta)(dt).
\]
There are finitely many jumps above each cutoff. Independence of Poisson
increments gives
$\mathbb E|Y_\epsilon-Y_\delta|^2
=\int_{\min(\epsilon,\delta)}^{\max(\epsilon,\delta)}e^{-t}\,dt$,
so $Y_\epsilon\to Y$ in $L^2$ as $\epsilon\downarrow0$. Conditioning
on the Poisson number of jumps gives
\[
   \mathbb E e^{zY_\epsilon}=
   \exp\!\left(z+\int_\epsilon^\infty
     (e^{-zt}-1+zt)t^{-2}e^{-t}\,dt\right)
       \qquad(\operatorname{Re}z>-1).
\]
The integral is absolutely convergent down to zero. For each such $z$
choose $r>1$ with $r\operatorname{Re}z>-1$. The same formula gives
$\sup_{0<\epsilon<1}\mathbb E|e^{zY_\epsilon}|^r<\infty$, permitting
passage to the limit by uniform integrability. The exponent at zero
cutoff has value $0$ and derivative $1$ at $z=0$. Its second
derivative is $\int_0^\infty e^{-(1+z)t}\,dt=1/(1+z)$, so it equals
$(1+z)\log(1+z)$. Taking $S=e^Y$ gives the claimed transform.

Let $X$ be $S$ times an independent mean-one exponential variable.
Its law has Mellin transform
$M(z)=\Gamma(1+z)(1+z)^{1+z}$ for $\operatorname{Re}z>-1$.
The moments $m_n=n!(n+1)^{n+1}$ satisfy
\[
  \frac{m_{n+1}/(2n+2)!}{m_n/(2n)!}
  =\frac{n+2}{2(2n+1)}\left(1+\frac1{n+1}\right)^{n+1}
  \longrightarrow \frac{e}{4}<1 .
\]
In particular $\mathbb E e^{\sqrt X}\le 2\mathbb E\cosh(\sqrt X)
=2\sum_{n\ge0}m_n/(2n)!<\infty$. If $(1/m_n)$ were a CZDS,
Lemma~\ref{lem:czdsmellin} would extend $1/M$ to an entire function
$E$. On $w>0$, analytic continuation would give
$E(w-1)=w e^{-w\log w}/\Gamma(1+w)$. This tends to zero when
$w\downarrow0$, so $G(w)=E(w-1)\Gamma(1+w)/w$ is holomorphic near
zero with $G(0)=1$. But on $w>0$,
$(G(w)-1)/w\sim-\log w$ as $w\downarrow0$, a contradiction.
\end{proof}

\subsection{An exact degree-72 counterexample certificate}
The specific sequence in Problem 37 admits a finite arithmetic check.
Write $m_j=j!/(j^2+j+1)$ and $L(x^j)=m_j$. In the basis $x^i/i!$,
the matrix of the quadratic form $u\mapsto L(u^2)$ in degree at most
36 has entries
\[
 B_{ij}=\binom{i+j}{i}\big/((i+j)^2+i+j+1),\qquad 0\le i,j\le36.
\]
Its first 36 leading principal minors are positive, and the last is
negative. This Python 3 certificate uses fraction-free elimination.
Each diagonal pivot before the final entry is the corresponding
leading principal determinant of the scaled integer matrix.
\begin{verbatim}
from math import comb, lcm

n = 37
den = [k*k + k + 1 for k in range(2*n-1)]
scale = lcm(*den)
h = [[comb(i+j, i)*(scale//den[i+j]) for j in range(n)]
     for i in range(n)]
prev = 1
for k in range(n-1):
    pivot = h[k][k]
    assert pivot > 0
    for i in range(k+1, n):
        for j in range(k+1, n):
            num = pivot*h[i][j] - h[i][k]*h[k][j]
            h[i][j], rem = divmod(num, prev)
            assert rem == 0
    prev = pivot
assert h[-1][-1] < 0
print("36 positive leading minors, then a negative determinant")
\end{verbatim}
For an explicit specification of the witness, let $B_0$ be the
leading $36$ by $36$ block of $B$ and take the vector $v_i=B_{i,36}$,
$0\le i<36$. Set $(c_0,\ldots,c_{35})^t=-B_0^{-1}v$, $c_{36}=1$,
and $u(x)=\sum_{i=0}^{36}c_ix^i/i!$. The Schur complement gives
$L(u^2)=\det B/\det B_0<0$. Put $\epsilon=-L(u^2)/2>0$ and
$f=T_m(u^2+\epsilon)$. Here $m_0=1$, so $f(1)<0$.
Its degree is 72 and its leading coefficient is positive, hence it has
at least two real zeros. But $T_\gamma f=u^2+\epsilon$ has all 72 zeros
nonreal. All coefficients in this specification are rational.

\section{Sharp constants for permuted gaps}\label{sec:garcia}
Rogelio Tom\'as Garc\'ia, \emph{A general lower bound for average
local discrepancy and an application to the Farey sequence},
\emph{Mathematics} 14 (2026), 2543, \S3, Conjecture~1,
\cite{tg} proposed the universal upper constant $4\sqrt{2/3}/9$ on empirical
grounds. The following theorem determines the optimal constant at
every fixed length and verifies the proposed universal value.

\medskip\noindent\textbf{Theorem.}
Let $0=a_0<a_1<\cdots<a_N=1$, let $g_j=a_j-a_{j-1}$, and put
$\sigma^2=N^{-1}\sum_j(g_j-1/N)^2$. For a uniform permutation
$\pi$, write $a_k^{\pi}=\sum_{j=1}^k g_{\pi(j)}$ and
$\overline r=\mathbb E\sum_{k=1}^N|a_k^\pi-k/N|$.
For $N\ge2$,
\[
 \overline r\ \le\ C_N\,\sigma N^{3/2},\qquad
 C_N=\frac{(N+1)\sqrt{N-1}}{3N^{3/2}}.
\]
The constant $C_N$ is attained for nonconstant gaps precisely when
all but one gap are equal. The smallest constant working for all $N$
is $C_3=4\sqrt{2/3}/9$. The case $N=1$ has $\overline r=0$.

\medskip\noindent\emph{Proof.}
Write $d_j=g_j-1/N$, $h_k=\sum_{j=1}^k d_{\pi(j)}$, and $h_0=h_N=0$.
Regard $h_0,\ldots,h_{N-1}$ as the vertices of a closed path. Each
cyclic shift of a uniform permutation remains uniform. Rerooting at
vertex $i$ subtracts $h_i$ from every vertex, hence
\[
 \frac{\overline r}{N}
 =\mathbb E\,\frac{1}{N^2}\sum_{i,j=0}^{N-1}|h_i-h_j|.
\]
This identity also applies to repeated gap values because all distinct
orders have equal numbers of labelled permutations.

For real numbers $t_1\le\cdots\le t_N$, write $\bar t$ for their mean
and $V=N^{-1}\sum_j(t_j-\bar t)^2$. The finite-sample
Gini mean difference satisfies
\[
 \frac1{N^2}\sum_{i,j}|t_i-t_j|
 =\frac2{N^2}\sum_{j=1}^N (2j-N-1)(t_j-\bar t)
 \le \frac2{\sqrt3}\sqrt{1-N^{-2}}\,\sqrt V .
\]
Indeed, the sum of the squared rank scores is $N(N^2-1)/3$.
Cauchy--Schwarz is an equality exactly for an arithmetic progression,
including the constant one.

Let $V_\pi$ be the empirical variance of the vertices. For two vertices
$k$ steps apart around the path, the increment is a size-$k$ sample
sum without replacement from the $d_j$. Its second moment is
$k(N-k)\|d\|_2^2/[N(N-1)]$, by summing the $k$ diagonal terms of
expectation $\|d\|_2^2/N$ and $k(k-1)$ off-diagonal terms of
expectation $-\|d\|_2^2/[N(N-1)]$. Thus
\[
 \mathbb EV_\pi
 =\frac1{2N^2}\sum_{i,j=0}^{N-1}\mathbb E(h_i-h_j)^2
 =\frac{\|d\|_2^2}{2N^2(N-1)}
       \sum_{k=0}^{N-1}k(N-k)
 =\frac{N+1}{12N}\|d\|_2^2.
\]
Applying the rank-score bound to the sorted vertices and using
$\mathbb E\sqrt{V_\pi}\le\sqrt{\mathbb E V_\pi}$ gives the claimed
$C_N$, since $\sigma N^{3/2}=N\|d\|_2$.

For increments $b,\ldots,b,-(N-1)b$, the vertices of any ordering
are a translated arithmetic progression of spacing $|b|$. Consequently
every ordering attains the rank-score bound and has the same variance.
Taking a sufficiently small nonzero $b$ gives strictly positive gaps.

For necessity, equality forces $V_\pi$ to be constant over all
permutations, and the vertices of each to form an arithmetic
progression after sorting. Suppose $N\ge4$ and fix two unequal
increments $u,v$ in the first two positions. Swapping them replaces
the vertex $u$ by $v$ and leaves the other $N-1$ vertices unchanged.
If $L$ is the sum of those common vertices, equality of the variances
gives $(N-1)(u+v)=2L$. Interchanging two successive increments in
the remaining $N-2$ positions changes $L$ by their difference, while
the left-hand side is fixed. Thus all increments outside any unequal
pair are equal. Fix such a pair with common outside value $b$.
If both members differed from $b$, the complement of the pair formed
from the first member and an outside $b$ would still contain the second
member and another $b$, a contradiction. For $N=3$ the three vertices
are $0,u,u+v$; one is the midpoint of the other two only when two of
$u,v,-u-v$ coincide. For $N=2$ the stated equality condition is automatic.
Finally, $9 C_N^2=1+x-x^2-x^3$ with $x=1/N$; its derivative
$(1+x)(1-3x)$ shows the universal maximum occurs at $N=3$.
\hfill$\square$


\subsection{The matching lower extremizers}
For completeness we give the lower bound from our preceding manuscript
and determine its equality cases. Write $R(S)=\overline r$, $a_j=g_j-1/N$
and $A=\|a\|_2$. The exact limits at fixed $N$ are
\[
 \frac{N+1}{3\sqrt2}A\ \le\ R(S)\ \le \frac{N+1}{3}\sqrt{1-1/N}\,A .
\]
For $0<p<1$, let $B_1(p),\ldots,B_N(p)$ be independent Bernoulli
variables of success probability $p$. With $T_p=\sum_j a_j(B_j(p)-p)$,
\begin{equation}
 \frac{R(S)}{N+1}=\int_0^1\mathbb E|T_p|\,dp .
 \label{eq:beta-perm}
\end{equation}
Indeed, the first $k$ entries of a uniform permutation form a uniform
$k$-subset. A fixed set $I$ of size $k$ is the set of Bernoulli
successes with probability $p^k(1-p)^{N-k}$. Its integral is
$1/((N+1)\binom Nk)$. Summing over sets gives~\eqref{eq:beta-perm},
since $\sum_j a_j=0$ and the two endpoint partial sums vanish.

\begin{lemma}[A sharp lower estimate at each sampling probability]
\label{lem:bias-bern}
For real $u_1,\ldots,u_N$, independent Bernoulli$(p)$ variables $B_j$
and $0\le p\le1$,
\[
 \mathbb E\left|\sum_j u_j(B_j-p)\right|
 \ge \sqrt2\,p(1-p)\left(\sum_j u_j^2\right)^{1/2}.
\]
Opposite coefficients of equal magnitude, with all other coefficients
zero, give equality for every $p$.
\end{lemma}

\begin{proof}
The endpoints and the zero vector are immediate. Write $q=1-p$ and
normalize $\sum_j u_j^2=1$, omitting zero terms.
If $|u_j|\ge 1/\sqrt2$ for some $j$, conditional Jensen applied
given $B_j$ bounds the left-hand side below by
$|u_j|\mathbb E|B_j-p|=2pq|u_j|$.

Suppose now that $|u_j|\le1/\sqrt2$ for all $j$. Put
$h(v)=1-2pq+2pq\cos v$. This is the characteristic function of the
difference of two independent Bernoulli$(p)$ variables. It is
nonnegative because $2pq\le1/2$. We use the following integral
inequality from Havrilla and Tkocz \cite[proof of Theorem~8,
Section~3.4, equation (23) in the preprint]{ht}. If
$h(v)=\rho_0+(1-\rho_0)\cos v$ with $1/2\le\rho_0\le1$, then
for every real $s\ge1$,
\[
 F_h(s):=\frac2\pi\int_0^\infty
   \bigl(1-h(v/\sqrt s)^s\bigr)\frac{dv}{v^2}
 \ge F_h(1)=1-\rho_0.
\]
Here the atom at zero is $\rho_0=1-2pq$.

Let $T=\sum_j u_j(B_j-p)$, with characteristic function $\phi_T$.
Independence and the weighted arithmetic--geometric mean inequality give
\[
 \operatorname{Re}\phi_T(v)
 \le \prod_j h(u_jv)^{1/2}
 \le \sum_j u_j^2 h(u_jv)^{1/(2u_j^2)} .
\]
The first absolute moment identity
$\mathbb E|T|=(2/\pi)\int_0^\infty(1-\operatorname{Re}\phi_T(v))\,dv/v^2$
follows from the corresponding cosine integral for $|x|$ and Tonelli's
theorem. Apply it to the preceding inequality, use evenness of $h$,
and substitute $v=\sqrt2\,w$. With $s_j=1/(2u_j^2)\ge1$ this yields
\[
 \mathbb E|T|\ge
 \frac1{\sqrt2}\sum_j u_j^2 F_h(s_j)
 \ge \frac{2pq}{\sqrt2}.
\]
The normalized vector $(1,-1,0,\ldots,0)/\sqrt2$ gives
$\mathbb E|B_1-B_2|/\sqrt2=\sqrt2\,pq$.
\end{proof}


Integrating the lemma gives the lower coefficient since
$\int_0^1p(1-p)\,dp=1/6$. For a nonconstant array, equality in this
integral requires pointwise equality for every $p$ by continuity.
Since $\sum a_j=0$, the outcomes with exactly one Bernoulli success give
$\mathbb E|T_p|=p\sum_j|a_j|+O_a(p^2)$ as $p\downarrow0$.
Equality therefore requires $\sum |a_j|=\sqrt2 A$.
The positive and negative coefficients each have total absolute mass
$L$, so $A^2\le 2L^2$, with equality exactly when each sign occurs once.
Thus the lower extremizers have coefficients $v,-v,0,\ldots,0$, up to
permutation. They attain equality by the lemma. Letting $N$ tend to infinity
gives the sharp dimension-independent lower constant $1/(3\sqrt2)$.

\subsection{How rarely a permutation matches the native Farey order}
Let $0=x_0<x_1<\cdots<x_N=1$ be the Farey points of order $n$; here
$N=\sum_{q\le n}\varphi(q)\sim 3n^2/\pi^2$. Set
$r_n=\sum_{k=1}^N|x_k-k/N|$. Let $r_{n,\pi}$ denote the same statistic
after a uniformly random permutation of the $N$ gaps. All probabilities
below refer to this uniform finite set. Write
$V(n)=(\log n)^{3/5}(\log\log n)^{-1/5}$ for $n>e^e$.

\medskip\noindent\textbf{Theorem.}
There are absolute $C,c>0$ for which, for every $n\ge2$ and $h\ge0$,
\[
 \mathbb P(r_{n,\pi}\le h)
 \le C\exp\{-c\min(n,n^2/(1+h)^2)\}.
\]
In particular, for all sufficiently large $n$, the choice $h=r_n$
gives $\mathbb P(r_{n,\pi}\le r_n)\ll\exp\{-\exp(c'V(n))\}$
unconditionally for some $c'>0$. Under ordinary RH, for each
$0<\eta<1$ this improves to $\ll\exp(-c_\eta n^{1-\eta})$.

\medskip\noindent\emph{Proof.}
The ordered adjacent denominator pairs are the coprime $q,r\le n$
with $q+r>n$, and have gap $1/(qr)$. This follows from the neighbouring
fraction determinant $1$ and the mediant denominator $q+r$.
Choose smaller gaps with $3n/4<q,r\le n$ and larger gaps with
$n/8<q\le n/4$, $7n/8<r\le n$. Their sizes differ by
at least $2/n^2$. Each set has at least $bN$ labels for an absolute
$b>0$ and large $n$, since each rectangle has $(6/\pi^2)$ times its
area plus $O(n\log n)$ coprime points by M\"obius inversion.

Take $m$ blocks of $2M$ positions each, with $M=\lfloor N/(2m)\rfloor$,
leaving surplus positions unused. Pair positions $j$ and $j+M$ inside
each block. Call a pair mixed if its labels come from opposite chosen
sets. A block has at least $b^2M$ mixed pairs except with probability
$O(e^{-c_0M})$. Indeed, its expected count is at least $2b^2M$.
When its $2M$ labels are exposed successively, different possible next
labels can be coupled by transposition, changing at most two pairs.
Azuma's inequality for the exposure martingale gives the claim.

Condition on unordered pair contents and all surplus labels. The
orientations of the pairs are independent fair signs. Denote by
$T_\ell$ the sum, over block $\ell$, of the centred partial sums of
the permuted gaps. Previous blocks enter only through their fixed
totals, so the $T_\ell$ are conditionally independent. Reversing a mixed
pair changes $T_\ell$ by at least $2M/n^2$. Thus its signed-sum
coefficient has absolute value at least $M/n^2\gg1/m$ when
$M\asymp n^2/m$. With $s\ge b^2M$ mixed pairs, conditional on all
other signs, an interval of length less than twice this coefficient
lower bound has probability at most
$2^{-s}\binom{s}{\lfloor s/2\rfloor}\ll M^{-1/2}$.
In fact, after making the coefficients positive, the subsets of
positive signs hitting the interval form an antichain (the
Erd\H{o}s--Littlewood--Offord argument \cite{erdoslo}). Covering the
interval $[-2h/m,2h/m]$ gives a bound $K(1+h)/\sqrt M$.

Choose $m=\lfloor\delta\min(n,n^2/(1+h)^2)\rfloor$ with $\delta$
a small positive absolute constant. If the quantity under the floor
is below two, the theorem follows by increasing $C$. Otherwise
$M\asymp n^2/m$, $M\gg m$ and $K(1+h)/\sqrt M\le1/16$.
Since $\sum_\ell |T_\ell|\le r_{n,\pi}$, at least half the blocks
must hit the last interval on the event $r_{n,\pi}\le h$.
Summing over choices of these blocks gives
$2^m(1/16)^{m/2}$ on the event all blocks have sufficiently many
mixed pairs. Adding $O(me^{-c_0M})$ proves the first bound.

For the unconditional specialization put
$D(x)=\#\{j\ge1:x_j\le x\}-Nx$, $M_0(y)=\sum_{k\le y}\mu(k)$ and
$\psi(t)=\{t\}-1/2$. M\"obius inversion gives
$\#\{j\ge1:x_j\le x\}=\sum_{d\le n}M_0(n/d)\lfloor dx\rfloor$
on $[0,1]$. Also $\sum_d M_0(n/d)=1$ and $\sum_d dM_0(n/d)=N$,
where $d\le n$. Consequently $D=-1/2-\sum_{d\le n}M_0(n/d)\psi(d\,\cdot)$.
The Fourier series of $\psi$ gives
$\int_0^1\psi(dx)\psi(ex)\,dx=(d,e)^2/(12de)$ by Parseval.
Use $M_0(y)\ll y\exp(-b_0V(y))$ for large $y$, from the classical
Vinogradov--Korobov region \cite{leongmertens}. For $d,e\le\sqrt n$
one has $V(n/d),V(n/e)\gg V(n)$. For the remaining terms bound
$|M_0(y)|\le y$. The kernel $(d,e)^2/(d^2e^2)$ has finite total sum
and its tail for $\max(d,e)>L$ is $O((1+\log L)/L)$: writing
$d=gu$, $e=gv$ with $g=(d,e)$ majorizes the tail by
$2\zeta(2)\sum_{gu>L}g^{-2}u^{-2}$, which has this bound.
Taking $L=\sqrt n$ in the formula for $\int D^2$ yields
$\|D\|_2\ll n e^{-b_1 V(n)}$ for some $b_1>0$.

Integrating the area between the empirical distribution and the
diagonal in the two coordinate directions gives
$\int_0^1|D(x)|\,dx=N\sum_j\int_{(j-1)/N}^{j/N}|x_j-t|\,dt$.
The triangle inequality on each such interval yields
$r_n\le \int_0^1|D(x)|\,dx+1/2\le\|D\|_2+1/2$.
The unconditional bound follows by substitution since $V(n)=o(\log n)$.
Under RH, Landau's criterion \cite{landau1924} gives
$r_n\ll_\eta n^{1/2+\eta/2}$. A second substitution proves the
claimed conditional rate.\hfill$\square$


\section*{Scope of the results}
The exact logarithmic densities in the Zhang race remain conditional
in the cited literature. In the polynomial classification we address
the complex-zero-decreasing parts of the cited questions; the ordinary
multiplier-sequence problems in parts (a) have different quantifiers.
For the native Farey order we remove RH from the conclusion that
matching its discrepancy is rare. The stronger rate
$\exp(-c_\eta n^{1-\eta})$ is proved here under RH.
The zero limit for the derivative thresholds leaves the value of
$\Lambda^{(0)}$ itself undetermined unconditionally.

\begin{thebibliography}{99}\small
\bibitem{KKL} H. Ki, Y.-O. Kim and J. Lee, On the de Bruijn--Newman
constant, \emph{Adv. Math.} \textbf{222} (2009), 281--306.
\href{https://doi.org/10.1016/j.aim.2009.04.003}{Publisher's text}.
\bibitem{LM} J. C. Lagarias and D. Montague, The integral of the Riemann
$\xi$-function, \emph{Comment. Math. Univ. St. Pauli} \textbf{60} (2011),
143--169. \href{https://arxiv.org/abs/1106.4348}{arXiv:1106.4348}.
\bibitem{Dobner} A. Dobner, A proof of Newman's conjecture for the extended
Selberg class, \emph{Acta Arith.} \textbf{201} (2021), 29--62.
\href{https://arxiv.org/abs/2005.05142v2}{Author's published-text version}, Theorem 4 and Lemma 3.
\bibitem{Ki22} H. Ki, On the de Bruijn--Newman constant, in
\emph{The Laguerre--P\'olya Class and Combinatorics}, Oberwolfach Report
13/2022, Theorem 1, p.~664.
\href{https://ems.press/content/serial-article-files/46952}{Workshop report}.
\bibitem{RT} B. Rodgers and T. Tao, The de Bruijn--Newman constant is
non-negative, \emph{Forum Math. Pi} \textbf{8} (2020), e6.
\href{https://arxiv.org/abs/1801.05914}{arXiv:1801.05914}.
\bibitem{cps} Y. Choie, M. Planat and P. Sol\'e, On Nicolas criterion
for the Riemann hypothesis, 2010, Conjecture 1 and Propositions 3--8.
\href{https://arxiv.org/abs/1012.3613}{arXiv:1012.3613v2}.
\bibitem{PT} D. J. Platt and T. S. Trudgian, On the first sign change
of $\theta(x)-x$, \emph{Math. Comp.} \textbf{85} (2016), 1539--1547.
\href{https://doi.org/10.1090/mcom/3021}{Publisher's text}, p.~1545.
\bibitem{oeis1} A. Karttunen, OEIS A342000, comments (2021).
\href{https://oeis.org/A342000}{Sequence entry}.
\bibitem{oeis5} A. Karttunen, OEIS A342455, comments (2021).
\href{https://oeis.org/A342455}{Sequence entry}.
\bibitem{rs62} J. B. Rosser and L. Schoenfeld, Approximate formulas for
some functions of prime numbers, \emph{Illinois J. Math.} \textbf{6}
(1962), 64--94. \href{https://doi.org/10.1215/ijm/1255631807}{Equation (3.16)}.
\bibitem{lpzh} J. D. Lichtman and C. Pomerance, The Erd\H{o}s conjecture
for primitive sets, \emph{Proc. Amer. Math. Soc. Ser. B} \textbf{6}
(2019), 1--14. \href{https://doi.org/10.1090/bproc/40}{Publisher's text}.
\bibitem{lmp} J. D. Lichtman, G. Martin and C. Pomerance, Primes in prime
number races, \emph{Proc. Amer. Math. Soc.} \textbf{147} (2019).
\href{https://doi.org/10.1090/proc/14569}{Publisher's text}.
\bibitem{dp} H. G. Diamond and J. Pintz, Oscillation of Mertens'
product formula, \emph{J. Th\'eor. Nombres Bordeaux} \textbf{21}
(2009), 523--533. \href{https://doi.org/10.5802/jtnb.687}{Publisher's text}.
\bibitem{CFbessel}
G. Csordas and T. Forg\'acs,
Multiplier sequences, classes of generalized Bessel functions and open
problems, \emph{J. Math. Anal. Appl.} \textbf{433} (2016), 1369--1389.
\href{https://doi.org/10.1016/j.jmaa.2015.08.047}{doi:10.1016/j.jmaa.2015.08.047}.
\href{https://arxiv.org/pdf/1505.00355v2}{arXiv:1505.00355v2} (2015).

\bibitem{CCsurveyZero}
T. Craven and G. Csordas, Composition theorems, multiplier sequences and
complex zero decreasing sequences, in \emph{Value Distribution Theory
and Related Topics}, Kluwer, 2004, 131--166.
\href{https://doi.org/10.1007/1-4020-7951-6_6}{doi:10.1007/1-4020-7951-6\_6}.

\bibitem{CCserdica1996}
T. Craven and G. Csordas, Problems and theorems in the theory of
multiplier sequences, \emph{Serdica Math. J.} \textbf{22} (1996), 515--524.

\bibitem{ISczds}
I. J. Schoenberg, On P\'olya frequency functions II. Variation-diminishing
integral operators of the convolution type, \emph{Acta Sci. Math.
(Szeged)} \textbf{12} (1950), 97--106.
\href{https://acta.bibl.u-szeged.hu/13626/}{Publisher's archive}.

\bibitem{BCCGczds}
A. Bakan, T. Craven, G. Csordas and A. Golub, Weakly increasing
zero-diminishing sequences, \emph{Serdica Math. J.} \textbf{22} (1996),
547--570. \href{https://www.math.bas.bg/serdica/1996/1996-547-570.pdf}
{Publisher's text}.

\bibitem{BRczds}
A. Bakan and S. Ruscheweyh, Solution of the Karlin problem for
zero-diminishing sequences satisfying a Carleman condition,
\emph{Proc. Amer. Math. Soc.} \textbf{136} (2008), 2665--2674.
\href{https://doi.org/10.1090/S0002-9939-08-09256-3}
{doi:10.1090/S0002-9939-08-09256-3}.
\bibitem{erdoslo}
P. Erd\H{o}s, On a lemma of Littlewood and Offord, \emph{Bull. Amer. Math. Soc.} \textbf{51} (1945), 898--902.
\bibitem{leongmertens}
E. S. Lee and N. Leong, New explicit bounds for Mertens function and the reciprocal of the Riemann zeta-function, arXiv:2208.06141.
\bibitem{tg}
R. Tom\'as Garc\'ia, A general lower bound for average local discrepancy and an application to the Farey sequence, \emph{Mathematics} \textbf{14} (2026), 2543. DOI 10.3390/math14142543.
\bibitem{ht}
A. Havrilla and T. Tkocz, Sharp Khinchin-type inequalities for symmetric discrete uniform random variables, \emph{Israel J. Math.} \textbf{246} (2021), 281--297. DOI 10.1007/s11856-021-2244-8. Preprint arXiv:1912.13345.
\bibitem{landau1924}
E. Landau, Bemerkungen zu der vorstehenden Abhandlung von Herrn Franel, \emph{Nachr. Ges. Wiss. G\"ottingen, Math.-Phys. Kl.} (1924), 202--206.

\end{thebibliography}
\end{document}
